% !TeX root = ../../main.tex
% =====================================================================
%  Análisis por elementos finitos con criterio
%  Responsable:
%  Última revisión:
% =====================================================================
\chapter{Análisis por elementos finitos con criterio}
\label{ch:elementos_finitos}

\begin{objetivos}
  \item Decidir si un problema merece un FEA o se resuelve más rápido y mejor a
        mano.
  \item Reconocer una singularidad numérica y no informar nunca de la tensión que
        aparece en ella.
  \item Justificar las condiciones de contorno de un análisis, que es de donde
        salen la mayoría de los errores.
  \item Hacer una comprobación manual que confirme o desmienta el resultado.
  \item Redactar un informe de análisis que otra persona pueda revisar.
\end{objetivos}

\fichasesion{90 minutos}{Capítulo \ref{ch:mecanica_aplicada}}%
{Ordenador con el software de simulación y una pieza real del coche}

Un programa de elementos finitos siempre da un resultado, y el resultado siempre
tiene el mismo aspecto convincente tanto si el modelo está bien planteado como si
no. La habilidad que hay que adquirir, por tanto, no es hacer correr el análisis,
sino \emph{saber cuándo el resultado no vale}. \textcite{cook_fea_2001} es la
referencia de fondo si se quiere entender lo que hay debajo.

% ---------------------------------------------------------------------
\section{Qué puede y qué no puede decirte un FEA}
\label{sec:elementos_finitos:01}

\begin{clave}
Un FEA no analiza tu pieza: analiza \textbf{el modelo que has construido de tu
pieza}. Si la geometría está simplificada, si las cargas son inventadas o si los
apoyos no se parecen a la realidad, el resultado es exacto \emph{para un problema
que no es el tuyo}.
\end{clave}

\begin{center}
\begin{tabularx}{\textwidth}{@{}XX@{}}
\toprule
\textbf{Bien, y de forma fiable} & \textbf{Mal, o solo con mucho cuidado} \\
\midrule
Rigidez y desplazamientos & Tensión máxima absoluta en una entalla \\
Comparar dos diseños entre sí & Vida a fatiga \\
Ver por dónde van las cargas & Uniones atornilladas y contactos \\
Decidir dónde quitar material & Cordones de soldadura \\
Modos y frecuencias propias & Pandeo, si solo se hace estático lineal \\
Reacciones en los apoyos & Materiales compuestos \\
\bottomrule
\end{tabularx}
\end{center}

Dos consecuencias prácticas. La primera: para \textbf{comparar} dos diseños, el
FEA es excelente aunque el modelo tenga defectos, siempre que los defectos sean
los mismos en ambos casos. La segunda: para dar un \textbf{número absoluto} que
va a justificar un coeficiente de seguridad, el listón es mucho más alto y hay
que verificar y validar.

Y una regla previa a todo lo demás: \textbf{si no sabes aproximadamente qué va a
salir, no lances el análisis}. Estima antes el orden de magnitud a mano; esa
estimación es lo que después te permitirá detectar que el resultado es absurdo.

% ---------------------------------------------------------------------
\section{Mallado y convergencia}
\label{sec:elementos_finitos:02}
\index{mallado}\index{convergencia}

\subsection{Elegir el tipo de elemento}

\begin{itemize}
  \item \textbf{Nunca uses tetraedros de primer orden para calcular tensiones.}
        Son excesivamente rígidos y dan resultados sistemáticamente bajos. Si el
        programa ofrece «borrador» o «rápido», normalmente es esto.
  \item \textbf{Tetraedros de segundo orden} son la opción por defecto razonable
        para piezas macizas.
  \item \textbf{Elementos lámina} para todo lo que sea de pared delgada: chapas,
        laminados, tubos si interesa el detalle local.
  \item \textbf{Elementos viga} para estructuras tubulares. Un análisis de rigidez
        torsional de un chasis tubular (capítulo \ref{ch:rigidez_torsional}) se
        hace con vigas: es más rápido, más limpio de interpretar y más fácil de
        comprobar a mano que un modelo sólido con cientos de miles de elementos.
\end{itemize}

\subsection{El estudio de convergencia}

Un resultado sin estudio de convergencia no es un resultado. El procedimiento es
mecánico:

\begin{enumerate}
  \item Elige la magnitud que te interesa: un desplazamiento, una tensión en un
        punto concreto, una reacción.
  \item Resuelve con tres o cuatro tamaños de malla decrecientes.
  \item Representa esa magnitud frente al número de grados de libertad.
  \item Si la curva se aplana, has convergido: quédate con la malla a partir de
        la cual el cambio es menor que un umbral que tú fijas (por ejemplo
        \SI{2}{\percent}).
  \item \textbf{Si la curva no se aplana y sigue creciendo, no es un problema de
        malla: es una singularidad.}
\end{enumerate}

\subsection{Singularidades}
\index{singularidad}

En una arista viva entrante, en una carga puntual o en un apoyo puntual, la
tensión teórica es infinita. Al refinar la malla, la tensión que informa el
programa crece sin parar y nunca converge. No es un fallo del programa: la
solución exacta del \emph{modelo} es infinita ahí.

\begin{ojo}
Informar de la tensión máxima que aparece en una arista viva es un error que un
juez detecta nada más ver la imagen. La respuesta correcta a «¿cuál es la tensión
máxima?» empieza por «esa zona es una singularidad; el valor que doy es el de
aquí, a esta distancia, y esta es la razón».
\end{ojo}

Qué hacer cuando aparece una:

\begin{itemize}
  \item \textbf{Modelar el radio de acuerdo real.} Casi siempre la pieza no tiene
        una arista viva: la tiene el modelo, porque alguien la simplificó.
  \item \textbf{Repartir la carga o el apoyo} en una superficie en lugar de en un
        punto o una arista.
  \item \textbf{Evaluar la tensión a una distancia} de la singularidad, indicando
        cuál, si el detalle local no es lo que se está estudiando.
  \item \textbf{Cambiar de magnitud}: si lo que importa es la rigidez, usa el
        desplazamiento, que converge mucho antes y no tiene este problema.
\end{itemize}

% ---------------------------------------------------------------------
\section{Condiciones de contorno}
\label{sec:elementos_finitos:03}
\index{condiciones de contorno}

Aquí es donde se pierden los análisis. La geometría suele estar bien y el
material también; lo que casi nunca se parece a la realidad es cómo se ha sujetado
y cargado la pieza.

\subsection{Sobrerrestricción}

Empotrar una cara entera porque «ahí va atornillada» añade una rigidez que no
existe y genera tensiones enormes justo en el borde del empotramiento, que además
son falsas. La pieza real está sujeta por unos tornillos que permiten cierto giro
y cierto deslizamiento.

\begin{center}
\begin{tabularx}{\textwidth}{@{}XX@{}}
\toprule
\textbf{En lugar de} & \textbf{Haz} \\
\midrule
Empotrar la cara de apoyo completa & Restringir solo lo necesario para evitar
movimiento de sólido rígido \\
Empotrar el agujero de un tornillo & Un punto remoto acoplado al agujero, o una
carga tipo apoyo distribuida \\
Fijar un nodo para «quitar el sólido rígido» & Alivio inercial, o tres
restricciones mínimas bien colocadas \\
Contacto pegado en toda la unión & Modelar el contacto real, o justificar por qué
pegado es conservador \\
\bottomrule
\end{tabularx}
\end{center}

\subsection{Simetría}

Solo se puede usar simetría cuando \textbf{son simétricas las tres cosas}:
geometría, cargas y apoyos. Un caso de carga en curva no es simétrico. Un caso de
frenada en recta sí puede serlo. Aprovechar la simetría cuando no la hay elimina
del modelo justo la parte de la carga que descompensa la pieza.

\subsection{Las cargas vienen de algún sitio}

Las fuerzas de un análisis no se eligen: salen de los casos de carga definidos en
el capítulo \ref{ch:casos_carga}, que a su vez salen de la dinámica del vehículo.
Un análisis con cargas «razonables» inventadas por quien lo hace no demuestra
nada, y en el evento de diseño la primera pregunta después de ver la imagen
siempre es de dónde sale ese número.

% ---------------------------------------------------------------------
\section{Verificación con cálculo manual}
\label{sec:elementos_finitos:04}
\index{verificación}

Antes de creerte un resultado, tres comprobaciones que cuestan poco:

\begin{enumerate}
  \item \textbf{Equilibrio.} La suma de las reacciones tiene que dar la suma de
        las cargas aplicadas, en las tres direcciones. Si no cuadra, hay una carga
        o un apoyo donde no debía.
  \item \textbf{Orden de magnitud a mano.} Simplifica la pieza a una viga, un
        tubo o una placa y aplica las fórmulas del capítulo
        \ref{ch:mecanica_aplicada}. Si el FEA y la mano difieren en un factor 2,
        puede ser la geometría; si difieren en un factor 10, hay un error.
  \item \textbf{Un caso conocido con la misma configuración.} Resuelve una viga
        en voladizo con tu mismo tipo de elemento, tu mismo mallado y tu mismo
        tipo de apoyo, y compárala con la solución analítica. Si eso no sale, lo
        demás tampoco.
\end{enumerate}

\begin{clave}
Conviene distinguir dos cosas que se confunden constantemente:
\textbf{verificación} es comprobar que estás resolviendo bien las ecuaciones
(malla, convergencia, equilibrio); \textbf{validación} es comprobar que las
ecuaciones que resuelves describen la realidad, y eso solo se hace midiendo la
pieza real. El banco de rigidez torsional del capítulo
\ref{ch:rigidez_torsional} es validación; un estudio de convergencia es
verificación. Hacen falta las dos.
\end{clave}

% ---------------------------------------------------------------------
\section{Cómo se documenta un análisis}
\label{sec:elementos_finitos:05}

Un análisis que no está documentado no se puede revisar, no se puede defender
ante un juez y no sirve el año que viene. El formato es una página, y siempre la
misma:

\begin{center}
\begin{tabularx}{\textwidth}{@{}lX@{}}
\toprule
\textbf{Apartado} & \textbf{Qué contiene} \\
\midrule
Objetivo & Qué decisión de diseño depende de este análisis \\
Geometría & Modelo usado y \emph{qué se ha simplificado y por qué} \\
Material & Designación y de dónde salen las propiedades \\
Caso de carga & Valores y referencia al documento del que salen \\
Contorno & Apoyos y cargas, con su justificación física \\
Malla & Tipo de elemento, tamaño y \textbf{estudio de convergencia} \\
Resultados & Magnitudes de interés, con leyenda y unidades \\
Comprobación & El cálculo manual y su comparación \\
Conclusión & Coeficiente de seguridad y \textbf{qué se cambia en el diseño} \\
\bottomrule
\end{tabularx}
\end{center}

Sobre las imágenes: llevan leyenda con escala y unidades, la escala no puede estar
autoajustada al valor de una singularidad, y si el desplazamiento está amplificado
hay que decir por cuánto. Sin esos tres datos, la imagen no se puede interpretar.

% ---------------------------------------------------------------------
%  Cierre del capítulo
% ---------------------------------------------------------------------

\begin{caso}[con qué calculamos]
El equipo está aprendiendo \textbf{Ansys}, y hasta ahora los análisis se han
hecho con los módulos de simulación integrados en \textbf{Fusion 360} y
\textbf{SolidWorks}. Conviene saber qué implica cada opción:

\begin{itemize}
  \item Los módulos integrados en el CAD son cómodos porque la geometría ya está
        ahí, pero dan poco control sobre el tipo de elemento, el mallado y los
        apoyos, y aplican valores por defecto que no siempre se ven. Sirven para
        comparar alternativas; cuesta más justificar con ellos un número absoluto.
  \item Ansys da ese control, y es la razón por la que merece la pena el esfuerzo
        de aprenderlo: permite elegir el orden del elemento y hacer de verdad el
        estudio de convergencia del apartado
        \ref{sec:elementos_finitos:02}.
\end{itemize}

Dos reglas mientras conviven las tres herramientas:

\begin{enumerate}
  \item \textbf{Cada persona resuelve el caso de referencia en la herramienta que
        vaya a usar} antes de analizar nada del coche. Es la única forma de saber
        qué hacen los valores por defecto de tu programa.
  \item \textbf{No se comparan resultados obtenidos con programas distintos} para
        decidir entre dos diseños: estarías comparando las herramientas, no los
        diseños. Una comparación se hace siempre con el mismo montaje.
\end{enumerate}

\redactar{Añadir aquí, cuando exista, un análisis real completo siguiendo la
plantilla del apartado \ref{sec:elementos_finitos:05}. Lo ideal es uno en el que
el resultado obligara a cambiar el diseño, y mejor todavía si se pudo contrastar
con una pieza rota o con una medida.}
\end{caso}

\begin{ojo}
\begin{itemize}
  \item Informar de la tensión máxima en una arista viva.
  \item Decir «ha convergido» refiriéndose a que el solver terminó, no a que la
        malla esté convergida. Son cosas distintas.
  \item Comparar dos diseños con mallas diferentes: la diferencia que ves puede
        ser solo la malla.
  \item Dar un coeficiente de seguridad contra un caso de carga que nadie ha
        definido por escrito.
  \item Comprobar el límite elástico en una pieza cuyo problema real es la
        rigidez o el pandeo.
  \item Mallar más fino para «afinar el resultado» sin comprobar si lo que hay
        debajo es una singularidad.
  \item Enseñar la imagen deformada con factor \num{200} sin decirlo.
\end{itemize}
\end{ojo}

\begin{entregable}
\begin{enumerate}
  \item La \textbf{plantilla de informe de análisis} del equipo, con los nueve
        apartados de la tabla anterior.
  \item Un \textbf{caso de referencia resuelto} (viga en voladizo) con la
        configuración habitual del equipo, contrastado con la solución analítica,
        guardado para que cualquiera pueda comprobar su montaje.
  \item El \textbf{checklist de revisión de un FEA}, que usa el revisor antes de
        dar por bueno un análisis ajeno.
\end{enumerate}
\end{entregable}

\begin{juez}
\begin{enumerate}
  \item ¿De dónde salen las cargas de este análisis?
  \item ¿Cómo has sujetado la pieza en el modelo y por qué se parece eso a la
        realidad?
  \item ¿Has hecho estudio de convergencia? Enséñamelo.
  \item Esa tensión máxima está en una arista viva. ¿Qué significa ese número?
  \item ¿Con qué has comprobado que este resultado es correcto?
  \item ¿Qué has cambiado en el diseño después de este análisis?
\end{enumerate}
\end{juez}
